\section{Problem Formulation}
Consider a leader indexed by $0$ and $N$ controlled trucks $i\in\{1,\ldots,N\}$ on a road coordinate. Each truck has physical state and brake commands
\begin{equation}
x_i=[p_i,v_i,b_i,r_i,T_i]^\top,\qquad
u_i=[u_i^f,u_i^a]^\top,
\end{equation}
where $p_i,v_i$ are longitudinal position and speed, $b_i,r_i$ are \emph{realized} friction and auxiliary forces, $T_i$ is friction-brake temperature, and $u_i^f,u_i^a$ are their distinct \emph{commands}. The hard QP therefore has exactly two command coordinates. Vehicle $i$ has parameter vector
\begin{equation}
\vartheta_i=(m_i,L_i,c_{r,i},C_{d,i}A_i,\tau_i^f,\tau_i^a,C_i,k_i,\eta_i,\bar b_i),
\end{equation}
where $m_i,L_i,c_{r,i},C_{d,i}A_i$ denote mass, length, rolling resistance, and drag area; $\tau_i^f,\tau_i^a$ are actuator time constants; and $C_i,k_i,\eta_i,\bar b_i$ denote lumped thermal capacity, cooling, friction heating, and cold friction-force capability. Tables~\ref{tab:parameters} and~\ref{tab:nomenclature} summarize provenance and notation.
Hard feasibility is not a function of $x_i$ alone because command memory under zero-order hold, gear availability, dwell logic, and communication validity depend on memory. At sample $k$ the controller therefore uses the augmented hybrid state
\begin{equation}\label{eq:chi}
\begin{aligned}
\chi_{i,k}=[&p_{i,k},v_{i,k},b_{i,k},r_{i,k},T_{i,k},
u^f_{i,k-1},u^a_{i,k-1},g_{i,k},\\
&\tau^g_{i,k},a^{\rm msg}_{i,k},d_{i,k},
v_{i-1,k},a_{i-1,k}]^\top .
\end{aligned}
\end{equation}
Here $g_{i,k}$ is the discrete gear state, $\tau^g_{i,k}$ is the elapsed dwell time in that gear state, $a^{\rm msg}_{i,k}$ is the age of the most recently accepted safety message, $d_{i,k}$ is the inter-vehicle bumper gap, and $v_{i-1,k}$ and $a_{i-1,k}$ are predecessor speed and acceleration. The last three quantities are required by the collision row and predecessor reachable tube. Every certificate below is consequently written as a function of $\chi_{i,k}$, not merely $x_i$.

The time-varying directed vehicle-to-vehicle (V2V) communication graph is $\mathcal G_t=(\mathcal V,\mathcal E_t)$. Agent $i$ observes
\begin{equation}
o_i=(v_i,b_i,r_i,T_i,u^f_{i,k-1},u^a_{i,k-1},g_i,\tau_i^g,
\vartheta_i,\hat\theta_i,d_i,\Delta v_i,
\mathcal M_i),
\end{equation}
where $d_i=p_{i-1}-p_i-L_{i-1}$, $\Delta v_i=v_{i-1}-v_i$, $\hat\theta_i$ is the locally estimated road grade, and $\mathcal M_i$ contains time-stamped permitted messages including the age needed for \eqref{eq:chi}. A message generated at $t_g$ and received at $t$ is replaced by a reachable tube $\mathcal R_{i-1}(t-t_g)$ formed from declared acceleration and jerk bounds. Packets inconsistent with local range/range-rate sensing are rejected, and the union of sensing and communication uncertainty is used in the robust margin. The critic uses global state only during training.

The policy must minimize tracking, comfort, braking-energy, coordination, and intervention costs while maintaining collision, temperature, fade, and actuator constraints. Jerk is a soft comfort constraint and therefore never determines the hard safety feasibility claim. Training, identification, validation, and held-out/out-of-distribution (OOD) test sets are disjoint.

At the transport level, a mission contains a descent profile $\theta(p)$, desired speed schedule, fleet order, payloads, ambient condition, and communication trace. We report
\begin{equation}
J_{\rm op}=\left[t_{\rm trip},E_{\rm fric},E_{\rm aux},e_v^{\rm rms},
e_d^{\rm rms},j^{\rm rms},n_{\rm stop}\right],
\end{equation}
where $E_{\rm fric}=\sum_i\int b_iv_i\,dt$ and $E_{\rm aux}=\sum_i\int r_iv_i\,dt$; $e_v^{\rm rms}$, $e_d^{\rm rms}$, and $j^{\rm rms}$ are speed-error, spacing-error, and jerk root-mean-square (RMS) quantities. Safety is reported separately from efficiency so that an unnecessarily slow method cannot appear superior merely because it avoids demanding maneuvers.

String stability is also separate from safety. With spacing error $e_i^d$, the empirical gains are
\begin{equation}
G_{2,i}=\frac{\|e_i^d\|_2}{\|e_{i-1}^d\|_2+\epsilon_g},\qquad
G_{\infty,i}=\frac{\|e_i^d\|_\infty}{\|e_{i-1}^d\|_\infty+\epsilon_g}.
\end{equation}
Here $\epsilon_g>0$ is a small denominator regularizer that prevents division by zero or a near-zero predecessor-error norm; it is not a safety parameter. We test whether $\max_iG_{2,i}\le1$ and $\max_iG_{\infty,i}\le1$ across platoon lengths; no analytical string-stability theorem is claimed.

\begin{table}[t]
\caption{Vehicle Parameters Required Per Class}
\label{tab:parameters}
\centering\normalsize
\begin{tabular}{@{}>{\raggedright\arraybackslash}p{0.22\columnwidth}>{\raggedright\arraybackslash}p{0.36\columnwidth}>{\raggedright\arraybackslash}p{0.24\columnwidth}@{}}
\toprule
Group & Parameters & Evidence \\
\midrule
Geometry & $m_i,L_i$, axle/load state & weighed vehicle or configuration \\
Resistance & $c_{r,i},C_{d,i}A_i$ & coast-down or source data \\
Actuator & $\tau_i^f,\tau_i^a,\bar u_i^f,\bar u_i^a(v,g)$ & friction and retarder step and power maps \\
Thermal & $C_i,k_i,\eta_i,T_{\rm crit,i}$ & dynamometer or descent fit \\
Fade & $\phi_i(T)$ and uncertainty & held-out force--temperature curve \\
Derivatives & $\dot\theta,\dot q^w,\dot F^r$, predecessor jerk bounds & held-out extrema or certified limits \\
Safety & $a_i^S,t_{h,i},d_{0,i}$ & validated lower bound/policy \\
\bottomrule
\end{tabular}
\end{table}

\begin{table}[t]
\caption{Core Symbols (SI Units Unless Stated)}
\label{tab:nomenclature}
\centering\normalsize
\begin{tabular}{@{}>{\raggedright\arraybackslash}p{0.31\columnwidth}>{\raggedright\arraybackslash}p{0.53\columnwidth}@{}}
\toprule
Symbol & Meaning \\
\midrule
$b,r$ & realized friction and auxiliary brake forces (N) \\
$u^f,u^a$ & commanded friction and auxiliary forces (N) \\
$T,T_a,T_{\rm crit}$ & brake, ambient, and critical temperatures (K) \\
$h^c,h^T,h^F,h^A$ & collision, thermal, fade, and auxiliary barriers \\
$A^f,A^a,D^c$ & collision-row coefficients and demand \\
$[L^f,U^f],[L^a,U^a]$ & effective command intervals after all individual rows \\
$M$ & exact conflict reserve in collision-row units \\
$g^{T0}$ & undivided low-speed state-only thermal reserve (K) \\
$\gamma,\mu$ & directional model and sampled-data margins \\
$H_{\rm QP}$ & safety-QP Hessian; $N_{\rm rec}$ is recovery count \\
\bottomrule
\end{tabular}
\end{table}


Figure~\ref{fig:coupled} summarizes the coupling among friction and auxiliary commands, their actuator lags, longitudinal dynamics, friction heating, and temperature-dependent fade.
\begin{figure*}[t]
\centering
\begin{tikzpicture}[font=\fontsize{9.4}{10.4}\selectfont,>=Latex,node distance=8mm,
 box/.style={draw,rounded corners,minimum width=25mm,minimum height=12mm,align=center,fill=gray!8}]
\node[box] (cmd) {admissible commands\\$u_i^f,u_i^a$};
\node[box,right=of cmd] (lag) {friction/auxiliary lags\\$\tau^f\dot b=-b+u^f$, $\tau^a\dot r=-r+u^a$};
\node[box,right=of lag] (long) {longitudinal plant\\$m_i\dot v_i=F_i^r-b_i-r_i$};
\node[box,below=12mm of long] (heat) {thermal plant\\$C_i\dot T_i=\eta_i b_i v_i-k_i(T_i-T_a)+q_i^w$};
\node[box,left=10mm of heat] (fade) {fade envelope\\$b_i\le \bar b_i\phi_i(T_i)$};
\draw[->] (cmd)--(lag);
\draw[->] (lag)--node[midway,above=5pt]{$b_i,r_i$}(long);
\draw[->] (lag.south) -- node[midway,left=5pt]{$b_i$} ++(0,-8mm) -| (fade.north);
\draw[->] ([xshift=8mm]lag.south) -- ++(0,-6mm) -| ([xshift=-10mm]heat.north);
\draw[->] (long)--node[midway,right=5pt]{$v_i$}(heat);
\draw[->] (heat)--node[midway,above=6pt]{$T_i$}(fade);
\draw[->] (fade.west) -- node[midway,above=6pt]{friction bound} (cmd.south |- fade.west) -- (cmd.south);
\end{tikzpicture}

\caption{Coupling among friction and auxiliary commands, their distinct actuator lags, longitudinal motion, friction heating, and the temperature-dependent fade envelope.}
\label{fig:coupled}
\end{figure*}
